% GENERATED FILE. Edit canonical lessons, then run npm run generate:books.
% canonical-source-hash: fnv1a64:c8a96799d7afaabb
% canonical-lessons: KA-C458-yatre, KA-C458-giridhama, KA-C458-karedoyyu, KA-C458-vapasagu, KA-C458-tangu

\chapter{Pilgrims and Hill Stations}
\label{ch:ka-pilgrims-and-hill-stations}
\hlchaptermodality{458}

\begin{chapteropening}
\textbf{By the end of this chapter:} \emph{I can say a pilgrimage, a hill station, to take along, to get back and to stay overnight.}

Complete the last lesson of chapter 458: I can say a pilgrimage, a hill station, to take along, to get back and to stay overnight.
\end{chapteropening}

\section[\texorpdfstring{\kn{ಯಾತ್ರೆ} (yātre)}{ಯಾತ್ರೆ (yātre)}]{\kn{ಯಾತ್ರೆ} (yātre) — a pilgrimage; a long journey}

\label{lesson:KA-C458-yatre}



\begin{quote}
Before the new one: say the Kannada for a spot, then the Kannada for to look splendid.
\end{quote}

\subsection*{You'll want to know: \kn{ಯಾತ್ರೆ}}
\textbf{\kn{ಯಾತ್ರೆ}} — \emph{yātre} — \textquotedblleft{}a pilgrimage; a long journey\textquotedblright{}.

\begin{quote}
\kn{ಅಜ್ಜಿ} \kn{ಪ್ರತಿ} \kn{ವರ್ಷ} \kn{ಯಾತ್ರೆಗೆ} \kn{ಹೋಗುತ್ತಾರೆ}.
\end{quote}

\begin{quote}
\emph{ajji prati varṣa yātrege hōguttāre.}
\end{quote}

\begin{quote}
Grandmother goes on a pilgrimage every year.
\end{quote}

\kn{ಯಾತ್ರೆ} alone can be any long journey; \kn{ಯಾತ್ರಿಕ} is a pilgrim.

\begin{cousinweb}
\textbf{\kn{ಯಾತ್ರೆ}} (\emph{yātre}) is borrowed from Sanskrit \emph{yātrā}, \textquotedblleft{}a journey, a festival procession\textquotedblright{} (CDIAL 10456), from \emph{yā}, \textquotedblleft{}to go\textquotedblright{}.
\end{cousinweb}

\begin{grammarlens}[title={What you've built}]
One more word for getting around.
\end{grammarlens}

\subsection*{Your turn}
\begin{itemize}
\raggedright
  \item \emph{Say it:} \emph{yātre}
  \item \emph{Say it:} \emph{yātre}, once more
  \item \emph{Say it:} \emph{ajji prati varṣa yātrege hōguttāre.} — Grandmother goes on a pilgrimage every year.
  \item \emph{Recall:} say the Kannada for a spot, then the Kannada for to look splendid, then say \emph{yātre} again
\end{itemize}

\begin{tcolorbox}[breakable,colback=teal!4,colframe=teal!35!black,title={Before you move on}]
What does \kn{ಯಾತ್ರೆ} mean? (\textquotedblleft{}A pilgrimage; a long journey\textquotedblright{}.) Say it once more.
\end{tcolorbox}
\section[\texorpdfstring{\kn{ಗಿರಿಧಾಮ} (giridhāma)}{ಗಿರಿಧಾಮ (giridhāma)}]{\kn{ಗಿರಿಧಾಮ} (giridhāma) — a hill station}

\label{lesson:KA-C458-giridhama}



\begin{quote}
Before the new one: say the Kannada for to look splendid, then the Kannada for a pilgrimage.
\end{quote}

\subsection*{You'll want to know: \kn{ಗಿರಿಧಾಮ}}
\textbf{\kn{ಗಿರಿಧಾಮ}} — \emph{giridhāma} — \textquotedblleft{}a hill station\textquotedblright{}.

\begin{quote}
\kn{ಬೇಸಿಗೆಯಲ್ಲಿ} \kn{ಗಿರಿಧಾಮದಲ್ಲಿ} \kn{ತಣ್ಣಗೆ} \kn{ಇರುತ್ತದೆ}.
\end{quote}

\begin{quote}
\emph{bēsigeyalli giridhāmadalli taṇṇage iruttade.}
\end{quote}

\begin{quote}
In summer it is cool at a hill station.
\end{quote}

\kn{ಗಿರಿ}, a hill, plus \kn{ಧಾಮ}, a dwelling.

\begin{cousinweb}
\textbf{\kn{ಗಿರಿಧಾಮ}} (\emph{giridhāma}) is a modern compound built from Sanskrit words: \emph{giri}, \textquotedblleft{}a rock, a mountain\textquotedblright{} (CDIAL 4161), and \emph{dhāman}, \textquotedblleft{}a dwelling place\textquotedblright{} (Monier-Williams).
\end{cousinweb}

\begin{grammarlens}[title={What you've built}]
One more word for getting around.
\end{grammarlens}

\subsection*{Your turn}
\begin{itemize}
\raggedright
  \item \emph{Say it:} \emph{giridhāma}
  \item \emph{Say it:} \emph{giridhāma}, once more
  \item \emph{Say it:} \emph{bēsigeyalli giridhāmadalli taṇṇage iruttade.} — In summer it is cool at a hill station.
  \item \emph{Recall:} say the Kannada for to look splendid, then the Kannada for a pilgrimage, then say \emph{giridhāma} again
\end{itemize}

\begin{tcolorbox}[breakable,colback=teal!4,colframe=teal!35!black,title={Before you move on}]
What does \kn{ಗಿರಿಧಾಮ} mean? (\textquotedblleft{}A hill station\textquotedblright{}.) Say it once more.
\end{tcolorbox}
\section[\texorpdfstring{\kn{ಕರೆದೊಯ್ಯು} (karedoyyu)}{ಕರೆದೊಯ್ಯು (karedoyyu)}]{\kn{ಕರೆದೊಯ್ಯು} (karedoyyu) — to take along; to escort}

\label{lesson:KA-C458-karedoyyu}



\begin{quote}
Before the new one: say the Kannada for a pilgrimage, then the Kannada for a hill station.
\end{quote}

\subsection*{You'll want to know: \kn{ಕರೆದೊಯ್ಯು}}
\textbf{\kn{ಕರೆದೊಯ್ಯು}} — \emph{karedoyyu} — \textquotedblleft{}to take along; to escort\textquotedblright{}.

\begin{quote}
\kn{ನಾನು} \kn{ಅಜ್ಜಿಯನ್ನು} \kn{ದೇವಸ್ಥಾನಕ್ಕೆ} \kn{ಕರೆದೊಯ್ದೆ}.
\end{quote}

\begin{quote}
\emph{nānu ajjiyannu dēvasthānakke karedoyde.}
\end{quote}

\begin{quote}
I took grandmother to the temple.
\end{quote}

Past (I …): \kn{ಕರೆದೊಯ್ದೆ}, \emph{karedoyde}.

\kn{ಕರೆ}, to call, plus \kn{ಒಯ್ಯು}, to take away: to call someone and take them.

\begin{cousinweb}
\textbf{\kn{ಕರೆದೊಯ್ಯು}} (\emph{karedoyyu}) is a native Dravidian compound: \emph{kare}, \textquotedblleft{}to call\textquotedblright{} (DEDR 1291), and \emph{oy}, \textquotedblleft{}to carry off, to take along\textquotedblright{} (DEDR 984).
\end{cousinweb}

\begin{grammarlens}[title={What you've built}]
One more word for getting around.
\end{grammarlens}

\subsection*{Your turn}
\begin{itemize}
\raggedright
  \item \emph{Say it:} \emph{karedoyyu}
  \item \emph{Say it:} \emph{karedoyyu}, once more
  \item \emph{Say it:} \emph{nānu ajjiyannu dēvasthānakke karedoyde.} — I took grandmother to the temple.
  \item \emph{Recall:} say the Kannada for a pilgrimage, then the Kannada for a hill station, then say \emph{karedoyyu} again
\end{itemize}

\begin{tcolorbox}[breakable,colback=teal!4,colframe=teal!35!black,title={Before you move on}]
What does \kn{ಕರೆದೊಯ್ಯು} mean? (\textquotedblleft{}To take along; to escort\textquotedblright{}.) Say it once more.
\end{tcolorbox}
\section[\texorpdfstring{\kn{ವಾಪಸಾಗು} (vāpasāgu)}{ವಾಪಸಾಗು (vāpasāgu)}]{\kn{ವಾಪಸಾಗು} (vāpasāgu) — to get back; to return}

\label{lesson:KA-C458-vapasagu}



\begin{quote}
Before the new one: say the Kannada for a hill station, then the Kannada for to take along.
\end{quote}

\subsection*{You'll want to know: \kn{ವಾಪಸಾಗು}}
\textbf{\kn{ವಾಪಸಾಗು}} — \emph{vāpasāgu} — \textquotedblleft{}to get back; to return\textquotedblright{}.

\begin{quote}
\kn{ಯಾತ್ರೆ} \kn{ಮುಗಿಸಿ} \kn{ನಾವು} \kn{ಮನೆಗೆ} \kn{ವಾಪಸಾದೆವು}.
\end{quote}

\begin{quote}
\emph{yātre mugisi nāvu manege vāpasādevu.}
\end{quote}

\begin{quote}
After the pilgrimage we got back home.
\end{quote}

Past (I …): \kn{ವಾಪಸಾದೆ}, \emph{vāpasāde}.

\kn{ವಾಪಸ್}, back, is common in speech: \kn{ವಾಪಸ್} \kn{ಕೊಡು}, to give back.

\begin{cousinweb}
\textbf{\kn{ವಾಪಸಾಗು}} (\emph{vāpasāgu}) joins \emph{vāpas}, \textquotedblleft{}back\textquotedblright{}, usually traced to Persian \emph{vāpas}, \textquotedblleft{}back, returned\textquotedblright{}, through Hindustani, and the native verb \kn{ಆಗು}, to become.
\end{cousinweb}

\begin{grammarlens}[title={What you've built}]
One more word for getting around.
\end{grammarlens}

\subsection*{Your turn}
\begin{itemize}
\raggedright
  \item \emph{Say it:} \emph{vāpasāgu}
  \item \emph{Say it:} \emph{vāpasāgu}, once more
  \item \emph{Say it:} \emph{yātre mugisi nāvu manege vāpasādevu.} — After the pilgrimage we got back home.
  \item \emph{Recall:} say the Kannada for a hill station, then the Kannada for to take along, then say \emph{vāpasāgu} again
\end{itemize}

\begin{tcolorbox}[breakable,colback=teal!4,colframe=teal!35!black,title={Before you move on}]
What does \kn{ವಾಪಸಾಗು} mean? (\textquotedblleft{}To get back; to return\textquotedblright{}.) Say it once more.
\end{tcolorbox}
\section[\texorpdfstring{\kn{ತಂಗು} (taṅgu)}{ತಂಗು (taṅgu)}]{\kn{ತಂಗು} (taṅgu) — to stay overnight; to halt}

\label{lesson:KA-C458-tangu}



\begin{quote}
Before the new one: say the Kannada for to take along, then the Kannada for to get back.
\end{quote}

\subsection*{You'll want to know: \kn{ತಂಗು}}
\textbf{\kn{ತಂಗು}} — \emph{taṅgu} — \textquotedblleft{}to stay overnight; to halt\textquotedblright{}.

\begin{quote}
\kn{ದಾರಿಯಲ್ಲಿ} \kn{ನಾವು} \kn{ಒಂದು} \kn{ರಾತ್ರಿ} \kn{ಹಳ್ಳಿಯಲ್ಲಿ} \kn{ತಂಗಿದೆವು}.
\end{quote}

\begin{quote}
\emph{dāriyalli nāvu ondu rātri haḷḷiyalli taṅgidevu.}
\end{quote}

\begin{quote}
On the way we stayed one night in a village.
\end{quote}

Past (I …): \kn{ತಂಗಿದೆ}, \emph{taṅgide}.

A stop on a journey; \kn{ಇರು} is to stay in general.

\begin{cousinweb}
\textbf{\kn{ತಂಗು}} (\emph{taṅgu}) is native Dravidian: DEDR 3014 lists Kannada \emph{taṅgu}, \textquotedblleft{}to stop, to stay, to sojourn, to lodge\textquotedblright{}.
\end{cousinweb}

\begin{grammarlens}[title={What you've built}]
One more word for getting around.
\end{grammarlens}

\subsection*{Your turn}
\begin{itemize}
\raggedright
  \item \emph{Say it:} \emph{taṅgu}
  \item \emph{Say it:} \emph{taṅgu}, once more
  \item \emph{Say it:} \emph{dāriyalli nāvu ondu rātri haḷḷiyalli taṅgidevu.} — On the way we stayed one night in a village.
  \item \emph{Recall:} say the Kannada for to take along, then the Kannada for to get back, then say \emph{taṅgu} again
\end{itemize}

\begin{tcolorbox}[breakable,colback=teal!4,colframe=teal!35!black,title={Before you move on}]
What does \kn{ತಂಗು} mean? (\textquotedblleft{}To stay overnight; to halt\textquotedblright{}.) Say it once more.
\end{tcolorbox}
